\documentclass[12pts]{article}

\usepackage[utf8]{inputenc}
\usepackage[spanish]{babel}

\usepackage{amssymb, amsmath}

\usepackage{geometry}
\geometry{margin=2.54cm}

\pagestyle{plain}

\begin{document}

	\begin{titlepage}
		\centering
		{\scshape\LARGE Universidad Aztlán \par}
		\vspace{1cm}
		{\scshape\Large Campus Playa del Carmen \par}
		\vspace{3cm}
		{\scshape\Huge Integrales \par}
		\vspace{1cm}
		{\itshape\Large Ejercicios \par}
		\vspace{3cm}
		{\Large José Rosendo Camal Teh \par}
		\vfill
		{\Large Ingeniería en Sistemas Computacionales \par}
		{\Large Segundo Cuatrimestre \par}
		{\Large Cálculo Integral \par}
		\vfill
		{\Large Profe. Xicontecatl Carreto \par}
		\vfill
		{\Large 26 de marzo de 2026 \par}
		\vfill
		{\Large Playa del Carmen, Quintana Roo \par}
	\end{titlepage}

	\section*{Cálculo Integral}
	\paragraph{\textbf{Instrucciones:}} Calcular las siguientes integrales inmediatas.

	% Desarrollo del ejercicio 1
	\subsection*{Ejercicio 1}

	\[ \int \sqrt{25x} \, dx \]

	Desarrollamos la integral:

	\begin{align*}
		\int \sqrt{25x} \, dx &= 5 \int \sqrt{x} \, dx \\
		&= 5 \cdot \frac{ x^{ \frac{3}{2} } }{ \frac{3}{2} } + C \\
		&= 5 \cdot \frac{2}{3} \cdot x^{ \frac{3}{2} } + C \\
		&= \frac{10}{3} x \sqrt{x} + C
	\end{align*}

	Comprobamos:

	\begin{align*}
		\sqrt{25x} &= \frac{d}{dx} \left( \frac{10}{3} x \sqrt{x} + C \right) \\
		&= \frac{d}{dx} \left( \frac{10}{3} x^{ \frac{3}{2} } + C \right) \\
		&= \frac{10}{3} \cdot \frac{3}{2} x^{ \frac{1}{2} } \\
		&= 5 \sqrt{x}
	\end{align*}

	Resultado:
	
	\[ \int \sqrt{25x} \, dx = \frac{10}{3} x \sqrt{x} + C \]

	\newpage
	% Desarrollo del ejercicio 2
	\subsection*{Ejercicio 2}

	\[ \int \frac{2}{ \sqrt{ x^{4} }} \, dx \]

	Desarrollamos la integral:

	\begin{align*}
		\int \frac{2}{ \sqrt{ x^{4} }} \, dx &= 2 \int x^{-2} \, dx \\
		&= -2x^{-1} + C
	\end{align*}

	Comprobamos:

	\begin{align*}
		\frac{2}{ \sqrt{ x^{4} }} &= \frac{d}{dx} \left( -2x^{-1} + C \right) \\
		&= 2x^{-2}
	\end{align*}

	Resultado:

	\[ \int \frac{2}{ \sqrt{ x^{4} }} \, dx = -2x^{-1} + C \]

	\newpage
	% Desarrollo del ejercicio 3
	\subsection*{Ejercicio 3}

	\[ \int \left( 3x^{3} + 5x^{2} - x \right) \, dx \]

	Desarrollamos la integral:

	\begin{align*}
		\int \left( 3x^{3} + 5x^{2} - x \right) \, dx &= \frac{3}{4} x^{4} + \frac{5}{3} x^{3} - \frac{1}{2} x^{2} + C
	\end{align*}

	Comprobamos:

	\begin{align*}
		3x^{3} + 5x^{2} - x &= \frac{d}{dx} \left( \frac{3}{4} x^{4} + \frac{5}{3} x^{3} - \frac{1}{2} x^2 + C \right) \\
		&= 3x^{3} + 5x^{2} - x
	\end{align*}

	Resultado:
	\[ \int \left( 3x^{3} + 5x^{2} - x \right) = \frac{3}{4} x^{4} + \frac{5}{3} x^{3} - \frac{1}{2} x^{2} + C \]

\end{document}
